\documentclass[11pt]{article}

\usepackage[margin=1in]{geometry}
\usepackage{times}
\usepackage{hyperref}
\usepackage{graphicx}
\usepackage{booktabs}
\usepackage{array}
\usepackage{listings}
\usepackage{xcolor}
\usepackage{enumitem}
\usepackage{amsmath}

\hypersetup{
  colorlinks=true,
  linkcolor=blue,
  urlcolor=blue,
  citecolor=blue,
  pdftitle={DataForge: A Streaming, Leak-Aware Pipeline for Synthetic LLM Fine-Tuning Datasets},
  pdfauthor={Ian Too}
}

\lstset{
  basicstyle=\ttfamily\small,
  backgroundcolor=\color{gray!8},
  frame=single,
  breaklines=true,
  columns=fullflexible,
  keywordstyle=\color{blue},
  commentstyle=\color{gray}
}

\title{DataForge: A Streaming, Leak-Aware Pipeline\\
for Synthetic LLM Fine-Tuning Datasets}
\author{Ian Too}
\date{2026}

\begin{document}
\maketitle

\begin{abstract}
Fine-tuning a language model on a narrow domain requires paired
instruction/response data that usually does not exist for that domain.
The common workaround --- crawl a site, chunk the text, and prompt a
stronger model to synthesize question--answer pairs from each chunk ---
is easy to prototype and easy to get subtly wrong. We describe
\textsc{DataForge}, an open-source pipeline that treats each of those
failure modes as a first-class design constraint rather than an
afterthought: HTTP failure semantics are honored (transient vs.\
permanent status codes, \texttt{Retry-After}, \texttt{robots.txt}
\texttt{Crawl-delay}); every artifact (page, chunk, sample) is
checkpointed so a multi-hour run is resumable rather than restartable;
collection and generation are overlapped through bounded, backpressured
queues to avoid idling an LLM behind a slow crawler; unattended LLM spend
is bounded by an explicit, concurrency-safe budget rather than only
indirect knobs like page count; and --- the concern motivating this
document --- dataset splitting is group-aware, so that near-duplicate
paraphrases generated from the same source page cannot appear on both
sides of a train/test boundary and silently inflate held-out evaluation
scores. We describe the system's architecture and the reasoning behind
each design decision, and evaluate it end to end on four websites: 831
samples for \$0.12, with page-disjoint splits on every site. The same
benchmark exposed six defects that unit tests had missed, including a
declared \texttt{Crawl-delay} that was not enforced while scraping; all
are fixed and reported. Of fourteen sites considered, ten were excluded
after a review of their terms and access policies, although
\texttt{robots.txt} alone would have permitted at least six of them. We also describe agent access through the Model Context
Protocol, situate the resulting synthetic, fine-tuning-oriented
dataset relative to retrieval-augmented generation as a complementary
route to reducing hallucination, discuss the responsible-use
considerations inherent to an unattended web-scraping and synthetic-data
tool, and situate the system relative to prior work in web-scale data
collection and synthetic instruction-tuning data generation.
\end{abstract}

\tableofcontents

\section{Motivation}

Instruction-tuning and domain adaptation of large language models (LLMs)
both benefit from data that reflects the target domain's vocabulary,
structure, and factual content \cite{wei2022finetuned, ouyang2022training}.
For many domains --- an organization's internal documentation, a
government agency's public guidance, a product's help center --- no
labeled instruction dataset exists, but a large corpus of prose does, in
the form of a website. \emph{Synthetic instruction generation} --- using
a capable model to read a passage and generate question--answer or
instruction--response pairs grounded in it --- has become a standard way
to bridge that gap \cite{wang2023selfinstruct, taori2023alpaca,
xu2023wizardlm}.

In practice, building this pipeline from scratch surfaces problems that
are easy to underestimate:

\begin{enumerate}[leftmargin=1.5em]
  \item \textbf{Web collection is not a solved sub-problem.} Real sites
    return a mix of transient errors (rate limiting, momentary 5xx),
    permanent errors (404, 403), and politeness constraints
    (\texttt{robots.txt} \texttt{Crawl-delay}) that a naive
    \texttt{requests.get} loop conflates.
  \item \textbf{Generation is the expensive, slow stage.} If collection
    and generation run strictly sequentially, the LLM --- the most
    expensive and highest-latency component --- sits idle while
    thousands of pages download.
  \item \textbf{A crawl-then-fail run is costly to repeat.} A pipeline
    that is not checkpointed forces a full re-crawl (and re-billing of
    every LLM call already made) after any interruption.
  \item \textbf{The resulting dataset violates the i.i.d.\ assumption
    that train/test splitting relies on.} Several synthetic samples are
    generated per chunk, and several chunks come from one page, so
    samples are not independent draws --- they cluster by source
    document.
\end{enumerate}

The fourth point is the one most likely to silently corrupt a result
rather than merely produce an error message, so we treat it in detail
in Section~\ref{sec:leakage}.

\section{System overview}

\begin{table}[h]
  \centering
  \small
  \caption{DataForge pipeline stages. A recipe (YAML) or the
  interactive wizard drives all six. Collection, Processing, and
  Generation run either strictly sequentially (batch mode) or as
  concurrent worker pools connected by bounded queues (streaming mode,
  Section~\ref{sec:streaming}). Quality scoring and export always run
  as batch stages because deduplication and group-aware splitting
  require visibility into the full sample set.}
  \label{tab:architecture}
  \begin{tabular}{@{}l p{2.1cm} p{2.1cm} l p{5.2cm}@{}}
    \toprule
    Stage & Consumes & Produces & Execution & Key mechanism \\
    \midrule
    Discovery  & Seed URLs & Candidate URLs & Batch &
      Sitemap and \texttt{robots.txt} parsing; bounded BFS fallback;
      headless retry for JS-rendered pages \\
    Collection & URLs & Pages (Markdown) & Batch or stream &
      Per-domain rate limit; transient vs.\ permanent status handling;
      \texttt{Crawl-delay} enforcement \\
    Processing & Pages & Token-bounded chunks & Batch or stream &
      Boilerplate stripping; overlapping chunks sized by
      \texttt{tiktoken} \\
    Generation & Chunks & Candidate samples & Batch or stream &
      LiteLLM provider portability; $n$ samples per chunk; shared
      cost/call budget \\
    Quality    & All samples & Approved samples & Batch &
      Heuristic and source-reference filters; exact and near-duplicate
      removal; fail-closed LLM judge \\
    Export     & Approved samples & JSONL / Parquet / CSV & Batch &
      Per-record lineage; group-aware train/validation/test split \\
    \bottomrule
  \end{tabular}
\end{table}

DataForge is organized as six stages (Table~\ref{tab:architecture}):
Discovery, Collection, Processing, Generation, Quality, and Export. A
run is specified either by a YAML \emph{recipe} --- a declarative,
version-controllable description of every stage's parameters --- or by
an interactive terminal wizard for exploratory use. The recipe format
is intentionally close to configuration-as-code practice in CI/CD
systems: an unknown key fails validation immediately with its exact
field path, rather than being silently ignored, and a recipe run
returns a distinct exit code for each class of outcome (invalid recipe,
zero URLs after filtering, paused, zero approved samples), so it
composes with a scheduler or CI pipeline without a human watching the
log.

\subsection{Discovery}
Candidate URLs are enumerated from \texttt{sitemap.xml} (including
sitemap indexes) and any \texttt{Sitemap:} directive in
\texttt{robots.txt}, before a single content page is fetched. Sites
without a usable sitemap fall back to a breadth-first crawl bounded by
depth and page count. Pages that look JavaScript-rendered (few raw
links relative to body size) are retried through a headless browser
(Playwright) when available. In interactive use, discovery is followed
by a review step --- substring, glob, or regex filtering plus per-URL
selection --- that a recipe replaces declaratively with
\texttt{source.include} / \texttt{source.exclude} / \texttt{source.max\_urls}.

\subsection{Collection}
An async HTTP client applies per-domain rate limiting and
\texttt{robots.txt} compliance. The central design decision here is
that \emph{not all non-200 responses mean the same thing}:

\begin{itemize}[leftmargin=1.5em]
  \item Transient statuses (\texttt{429}, \texttt{500}, \texttt{502},
    \texttt{503}, \texttt{504}) are retried with exponential backoff,
    honoring a server-supplied \texttt{Retry-After} header
    (delta-seconds or HTTP-date) capped at 120 seconds so a single slow
    URL cannot stall the run.
  \item Permanent statuses (\texttt{404}, \texttt{403}) are accepted as
    final and not retried, preserving crawl budget.
  \item A declared \texttt{robots.txt} \texttt{Crawl-delay} lowers the
    effective rate limit automatically; it can only make a crawl
    \emph{slower} than the configured rate, never faster, so an
    operator's explicit rate-limit ceiling is never silently overridden
    upward. The delay applies for the whole run, across every stage that
    makes requests; before version~2.4.0 it did not, and the benchmark of
    Section~\ref{sec:eval-defects} is what showed it.
\end{itemize}

A page is marked scraped only after a successful fetch, so a resumed
session retries exactly the URLs that previously failed, and none that
already succeeded.

\subsection{Processing}
Retrieved HTML is converted to Markdown, boilerplate (navigation,
footers, cookie notices) is stripped, and the remaining text is split
into overlapping, token-bounded chunks sized against the target model's
tokenizer (via \texttt{tiktoken}). Chunk size and overlap are
recipe-configurable; overlap exists so that a fact spanning a chunk
boundary is not lost to the generation stage entirely.

\subsection{Generation}
Each chunk is passed to an LLM (via LiteLLM, giving provider portability
across OpenAI, Anthropic, Google, Groq, Together AI, or a local Ollama
model) with a format-specific prompt template --- \texttt{qa},
\texttt{instruction}, \texttt{conversation}, or a user-supplied custom
system prompt --- and a recipe-level natural-language \texttt{goal}
string that steers topical relevance without hand-writing a full prompt
per domain. $n$ samples are requested per chunk (recipe:
\texttt{generation.n\_per\_chunk}), which is the direct source of the
non-independence discussed in Section~\ref{sec:leakage}.

\subsubsection{Cost control}
\label{sec:budget}
Because the pipeline is designed to run unattended (\texttt{dataforge
run}, typically cron'd), and because the quality judge (below) issues a
second LLM call per sample --- roughly doubling generation cost on its
own --- an operator's only cost control was, until recently, indirect:
bounding the number of pages or chunks fed into generation, fixed before
the provider or its pricing is even known. DataForge now exposes an
explicit, recipe-level budget --- \texttt{generation.max\_cost\_usd} and
\texttt{generation.max\_llm\_calls} --- enforced by a single tracker
shared across the generation and judge stages (and the streaming
pipeline's concurrent workers), so the cap bounds \emph{combined} spend
rather than either stage alone. The check happens before a network call
is dispatched, not after: once exhausted, a chunk is skipped cleanly
(the run finishes and reports how many calls were skipped) rather than
the process crashing or continuing to spend past the configured limit.
This mirrors the existing pattern of \texttt{crawl.max\_pages} bounding
collection, extended to the stage that actually incurs LLM cost.

\subsection{Quality}
Samples pass through layered filters of increasing cost, evaluated in
this order so that the LLM judge (the expensive step) never runs on
samples that a free filter would reject:

\begin{enumerate}[leftmargin=1.5em]
  \item A heuristic pre-filter on answer/question length and refusal
    phrases.
  \item A source-reference filter that rejects samples referring to
    ``the document'' or ``the passage'' rather than the subject matter
    --- a model trained on such samples learns to cite a source it will
    never have access to at inference time.
  \item Deduplication across the full session (Section~\ref{sec:dedup}).
  \item An LLM judge that reads each sample alongside its source chunk
    and scores it 1--5 on groundedness, rejecting anything below a
    configurable threshold. This step \emph{fails closed}: a sample the
    judge cannot score (malformed output, API error, missing key, or a
    call skipped by the budget cap of Section~\ref{sec:budget}) is
    rejected rather than passed through by default.
\end{enumerate}

The judge is empirically reliable on unambiguous failures (contradicted
facts, unsupported claims, source self-reference) and less reliable near
the threshold on answers that are vague but not technically wrong ---
consistent with known limitations of LLM-as-judge evaluation
\cite{zheng2023judging}. Operators who prefer precision over recall can
raise the score threshold.

\subsubsection{Deduplication: exact and near-duplicate}
\label{sec:dedup}
An earlier implementation fingerprinted a sample by MD5-hashing only the
first 200 characters of its concatenated message content, which is
wrong in both directions: two unrelated samples that happen to share an
opening sentence (common in templated agency guidance) collide as
false-positive duplicates, while the near-paraphrases that
\texttt{generation.n\_per\_chunk}$>1$ produces by design --- differing
in wording beyond the first 200 characters, which is the normal, expected
shape of asking an LLM for several question--answer pairs from the same
passage --- are never caught, since they hash differently. The exact-hash
check now covers the full message content, fixing the false-positive
direction, and is supplemented with a near-duplicate pass: token-set
Jaccard similarity computed only against other samples generated from the
\emph{same source chunk} (bounded comparison set, avoiding an $O(n^2)$
scan over the whole session), rejecting anything at or above a
configurable threshold (\texttt{quality.near\_dup\_threshold}, default
0.85). Without this, a session could log ``deduplication: passing'' while
the surviving samples were still substantially redundant paraphrases of
one another --- distinct from, and prior to, the cross-page leakage that
group-aware splitting (Section~\ref{sec:leakage}) addresses.

\subsection{Export}
Approved samples are written as JSONL, Parquet, and CSV, and optionally
published to the HuggingFace Hub or Kaggle. Every record carries
\texttt{page\_id}, \texttt{chunk\_id}, \texttt{chunk\_index}, and
\texttt{source\_url} regardless of whether a split is requested, because
that lineage is what makes a correct split possible even after the
fact. This lineage is not currently preserved uniformly across every
export format, however: the Unsloth/ShareGPT-formatted companion export
(\texttt{dataset\_unsloth.jsonl}) discards \texttt{page\_id},
\texttt{chunk\_id}, \texttt{chunk\_index}, and \texttt{source\_url},
retaining only the \texttt{conversations} array that loader format
expects. A consumer training directly from that file specifically ---
as opposed to the plain JSONL/Parquet/CSV exports, which retain full
lineage --- loses the ability to cite a generated sample's source
document, which matters wherever attribution or provenance is required
downstream (Section~\ref{sec:limitations}, project issue \#19).

\subsubsection{Dataset introspection}
Because quality scoring, deduplication, and splitting all happen inside
a single run, their outcomes were, until recently, visible only as
transient log lines during that run --- not queryable afterward, and not
available at all for a session viewed later. A \texttt{dataforge stats
<session\_id>} command now reports, read-only, over data the pipeline
already computes and now persists (the per-sample rejection reason is
written to the session database rather than only logged): approval
rate, a rejection-reason breakdown, the quality-score distribution,
question/answer length statistics, and the realized train/validation/test
proportions from the most recent export. Realized proportions are
reported rather than compared against the recipe's configured target
ratios, since target ratios are not yet persisted per-session outside of
the export step itself; target-vs-realized comparison is a natural
follow-up once that is stored.

\section{Concurrency: streaming vs.\ batch}
\label{sec:streaming}

Run strictly sequentially, the four collection-through-generation stages
leave the LLM --- the highest-latency, highest-cost component --- idle
for the entire duration of the crawl. DataForge's streaming mode
(\texttt{stream: true}, the recipe default) instead runs collection,
chunking, and generation as three concurrent worker pools connected by
bounded queues:

\[
\text{urls} \;\longrightarrow\; [\text{scrape pool}] \;\longrightarrow\; \text{pages} \;\longrightarrow\; [\text{chunk pool}] \;\longrightarrow\; \text{chunks} \;\longrightarrow\; [\text{LLM pool}] \;\longrightarrow\; \text{samples}
\]

Pool sizes are tuned independently, since each stage is bound by a
different resource: scraping is rate-limit bound (typically one or a few
requests per second per domain), while generation is bound by LLM
provider concurrency and cost. Queues between pools are bounded
(\texttt{DATAFORGE\_STREAM\_QUEUE\_SIZE}) to apply backpressure, so a
crawler that briefly outpaces the LLM cannot exhaust memory buffering
unprocessed pages. A failure at any single stage --- an unparsable page,
a malformed LLM response --- is isolated, logged, and skipped without
halting the pool; if generation fails fatally (invalid API key,
provider outage), scraping and chunking are allowed to finish so that a
subsequent resume only has to re-run generation. Because every unit of
work (URL, page, chunk) is checkpointed independently, resuming a
streaming session replays exactly the unfinished work, and quality
scoring and export remain deliberately batch-only: deduplication and
group-aware splitting are properties of the \emph{whole} sample set and
cannot be evaluated on a stream in isolation.

Sessions are interchangeable between modes at resume time: a run paused
under batch execution can be resumed with streaming enabled, and vice
versa, since both modes checkpoint to the same schema.

\section{Leak-aware dataset splitting}
\label{sec:leakage}

\subsection{The problem}
A synthetic dataset produced this way has a clustered generative
structure: $n$ samples are drawn per chunk, and multiple chunks are
drawn per page. Two samples generated from the same chunk --- or from
adjacent, overlapping chunks of the same page --- are frequently
near-paraphrases of one another, both in question and in the grounding
text available to answer it. This is precisely the condition under
which a uniformly random train/test split violates the independence
assumption implicit in held-out evaluation: memorization of a
training-set paraphrase can trivially answer a test-set item drawn from
the same underlying passage, producing an evaluation score that
overstates generalization. This is a specific instance of the more
general \emph{data leakage} problem well documented in the machine
learning literature \cite{kaufman2012leakage}, and closely related to
the near-duplicate contamination concerns raised for pretraining corpora
\cite{lee2022deduplicating}.

\subsection{The mitigation}
DataForge's export stage implements \emph{group-aware} splitting: the
grouping key (\texttt{export.split.group\_by}, default \texttt{page})
determines the unit that must not span a split boundary. Groups ---
whole pages by default, or URLs or chunks --- are assigned to
train/validation/test as atomic units using a seeded random assignment
(\texttt{export.split.seed}) targeting the configured proportions, and
the resulting assignment is verified (no group's members appear in more
than one split) before any file is written, rather than assumed correct
from the assignment procedure alone. This is a direct application of
group-based cross-validation splitting (as in, e.g.,
\texttt{GroupKFold} / \texttt{GroupShuffleSplit} in
scikit-learn~\cite{pedregosa2011scikit}) to synthetic
instruction-dataset construction, where the natural grouping key is the
source document rather than a subject or session identifier.

\subsection{Why this matters in practice}
Absent this mechanism, a practitioner who fine-tunes and evaluates on a
randomly split version of the same dataset is liable to observe an
evaluation score that reflects memorization of near-duplicate training
content rather than generalization to unseen material --- a failure
mode that is invisible from the split code alone and only surfaces (if
at all) as an unexpectedly large drop in performance on truly
out-of-distribution data later. Because DataForge always emits source
lineage (\texttt{page\_id}, \texttt{chunk\_id}, \texttt{source\_url})
even when no split is requested at export time, a correct group-aware
split remains possible downstream even for datasets exported before
this feature was used.

\subsection{Row-level randomization within a split}
\label{sec:row-shuffle}
Group-aware splitting randomizes which group is assigned to which split
(Section~\ref{sec:leakage}, above), which is the property
that prevents leakage. It does not, by itself, randomize row
\emph{order} within a split's output file: groups are assigned
largest-first so that realized proportions track the configured targets
closely (variance in achievable proportions grows as group sizes become
uneven, and assigning the largest, most constraining groups while every
split is still empty minimizes it), and each group's members are then
written as a contiguous block. The resulting \texttt{dataset\_train.*}
file therefore contains long runs of samples from the same source page,
ordered by descending page size, rather than a random interleaving of
rows.

Whether this matters depends on the consumer. Frameworks that shuffle
examples before or during training (HuggingFace \texttt{datasets},
Axolotl, TRL, in their default configurations) absorb this ordering
transparently. A consumer that reads the file sequentially, batches
without shuffling, or trains in a streaming fashion over the raw file
order is exposed to it directly: early batches within an epoch would be
dominated by whichever page happened to be largest, which can bias
early-training gradient statistics and any diagnostic computed
batch-by-batch rather than over a full epoch. This is a distinct concern
from group assignment and admits an independent fix: a deterministic,
seeded shuffle of each split's row list --- separate from, and applied
after, group assignment --- immediately before the file is written,
reusing \texttt{export.split.seed} for reproducibility. Group
boundaries are unaffected by this shuffle; only row order within a
split changes. We considered, and rejected, generating samples
independently per target split rather than splitting a single generated
pool: that would require deciding a sample's split membership
\emph{before} generation, which inverts the group-aware design above,
since the grouping key (the source page) is only well-defined once
collection has already happened, and would either duplicate LLM calls
per page across splits or simply recompute the same page-to-split
assignment DataForge already makes, at additional cost and no benefit
over shuffling the already-split result. At the time of writing this
row-level shuffle is not yet implemented; it is tracked as an open item
(Section~\ref{sec:limitations}, project issue \#20) rather than
presented as shipped.

\section{Empirical evaluation}
\label{sec:eval}

The preceding sections describe what DataForge is designed to do. This
section reports what it did when run end to end against real websites:
what it fetched, generated, approved and rejected, what that cost and how
long it took, whether the split leaked pages, and which of the design
claims above did not hold until the benchmark exposed them. It does not
measure the quality of a model fine-tuned on the output, and it does not
yet validate the LLM judge against human labels
(Section~\ref{sec:eval-limits}).

\subsection{Setup}
\label{sec:eval-setup}
We selected sites that differ in structure and domain, then crawled only
those that passed a terms review (Section~\ref{sec:eval-access}). Four
remained:
\begin{itemize}[leftmargin=1.5em]
  \item \textbf{Ready.gov} (U.S. disaster preparedness): a flat sitemap
    of 372 URLs, no \texttt{Crawl-delay}.
  \item \textbf{USCIS} (U.S. immigration): a sitemap index expanding to
    19{,}623 URLs, with \texttt{Crawl-delay: 10}.
  \item \textbf{The Python tutorial} (technical documentation): the
    site's sitemap lists only eight version landing pages, so discovery
    must fall back to a breadth-first crawl from the seed
    \texttt{/3/tutorial/}.
  \item \textbf{iantoo.space} (a small personal site, 8 sitemap URLs),
    owned by the author and crawled with the owner's consent.
\end{itemize}
Every run used the same recipe settings apart from scoping: at most 20
URLs (\texttt{source.max\_urls}), a \$1.00 spending cap
(\texttt{generation.max\_cost\_usd}), three samples per chunk, 512-token chunks with 64-token overlap, the LLM
judge enabled at \texttt{min\_judge\_score: 4}, streaming mode, and an
80/10/10 split grouped by page with seed 42. Generation and judging both
used \texttt{gpt-4o-mini} through LiteLLM. Runs were made on 2026-09-22 on
Windows~11 with Python~3.11, from commit \texttt{73593bf} (the first three
sites) and \texttt{4cc258b} (iantoo.space; it differs only in TLS
trust-store handling and the benchmark's own preflight check). The recipes,
the site list with its terms reviews, and the scripts that produced every
number below are in the repository under \texttt{evals/pipeline/}; each
run also writes a \texttt{run\_summary.json} (timing per stage, LLM calls,
cost, budget, errors) so the figures can be recomputed from disk.

\subsection{Access and terms}
\label{sec:eval-access}
The benchmark runner refuses to crawl a site until a person has recorded a
terms review, and runs a preflight that fetches \texttt{robots.txt} and the
seed page with DataForge's own User-Agent and skips any site that serves a
bot challenge or redirects to another host. Fourteen sites were considered;
Table~\ref{tab:access} summarizes the outcome.

\begin{table}[h]
  \centering
  \small
  \caption{Sites considered for the benchmark and why each was included or
  excluded. Quotations are from each site's own terms or policy pages,
  retrieved 2026-09-22.}
  \label{tab:access}
  \begin{tabular}{@{}>{\raggedright\arraybackslash}p{3.1cm}
                     >{\raggedright\arraybackslash}p{2.4cm}
                     >{\raggedright\arraybackslash}p{8.6cm}@{}}
    \toprule
    Site & Outcome & Basis \\
    \midrule
    USCIS & Included & ``Information presented on this WWW site is
      considered public information and may be distributed or copied.'' \\
    Ready.gov (FEMA) & Included & FEMA: ``Most material on FEMA.gov is free
      of copyright and may be copied and distributed without permission'';
      the page lists Ready.gov as a FEMA domain but does not extend these
      terms to it explicitly. \\
    Python docs & Included & Documentation under the PSF License Version~2. \\
    iantoo.space & Included & Owner's consent. \\
    \midrule
    KRA (Kenya) & Excluded: prohibited & Terms restrict ``data mining, data
      harvesting, data extracting or any other similar activity''; yet
      \texttt{robots.txt} returns 404, which by convention permits crawling. \\
    eCFR & Excluded: blocked & \texttt{robots.txt} allows the page, but
      requests are redirected to a bot block stating that ``programmatic
      access \ldots\ is limited to access to our extensive developer APIs.'' \\
    Kenya Power & Excluded: no grant & \texttt{robots.txt}: ``Allow all
      crawlers''; footer: ``All rights reserved.'' \\
    KEBS, Ministry of Health, Univ.\ of Nairobi, JKUAT, Strathmore (Kenya)
      & Excluded: no grant & Copyright or ``All Rights Reserved'' notices;
      no reuse statement found. \\
    FAA & Excluded: unclear & No reuse statement; the only policy text is a
      government-system banner: ``provided for U.S. Government-authorized
      use only.'' \\
    CDC & Excluded: not reviewed & \texttt{robots.txt} returned 403 to a
      generic User-Agent but 200 to DataForge's minutes later. \\
    \bottomrule
  \end{tabular}
\end{table}

Two observations follow. First, \texttt{robots.txt} and the terms that
actually govern reuse disagreed in both directions: KRA has no
\texttt{robots.txt} at all but prohibits data extraction in its terms;
eCFR's \texttt{robots.txt} permits the page while the site blocks
automated clients and points them to an API; and Kenya Power explicitly
invites all crawlers while reserving all rights in the content. A crawler
that treats \texttt{robots.txt} compliance as permission would have
crawled all three. Second, the outcome tracked jurisdiction. The U.S.
federal sites publish text that is public domain by statute
\cite{usc17sec105} and, in two cases, say so explicitly; every Kenyan site
we examined reserved its rights, and Kenyan law has no public-domain
exception for government works: copyright in works commissioned by the
Government vests in the Government, and works made by employees in the
course of their employment pass to the employer
(ss.~25(1), 31(1)(b) and 31(2) of \cite{kenyacopyright2001}). A dataset built only from sites
that grant reuse is therefore geographically skewed by construction, which
matters for any domain, such as public services, where the relevant
guidance is national.

\subsection{Pipeline results}
\label{sec:eval-results}

\begin{table}[h]
  \centering
  \small
  \caption{Benchmark results per site (each capped at 20 URLs and
  \$1.00). ``Approved'' is the LLM judge's verdict, not a human one; see
  Section~\ref{sec:eval-limits}.}
  \label{tab:bench}
  \begin{tabular}{@{}l r r r r r r r@{}}
    \toprule
    Site & Pages & Chunks & Samples & Approved & LLM calls & Cost (USD) & Wall (s) \\
    \midrule
    Ready.gov       & 20 &  64 & 192 & 189 (98.4\%) & 128 & 0.0248 &  91 \\
    USCIS           & 20 &  53 & 159 & 152 (95.6\%) & 106 & 0.0222 & 296 \\
    Python tutorial & 17 & 153 & 456 & 447 (98.0\%) & 305 & 0.0663 & 258 \\
    iantoo.space    &  6 &   8 &  24 &  23 (95.8\%) &  16 & 0.0031 &  14 \\
    \midrule
    Total           & 63 & 278 & 831 & 811 (97.6\%) & 555 & 0.1164 & 659 \\
    \bottomrule
  \end{tabular}
\end{table}

Table~\ref{tab:bench} gives the per-site results. All four runs completed
with no errors and no budget-skipped calls; the most expensive run used
6.6\% of its \$1.00 cap. Cost per approved sample was between
\$0.00013 and \$0.00015 on every site, so cost scaled with the amount of
text rather than with the site: the Python tutorial's pages are long (a
median of 2{,}152 words, against 831 for Ready.gov and 576 for USCIS),
giving 153 chunks from 17 pages. Every page was fetched successfully on
every site. On the Python tutorial the crawl fallback reached 40 pages, of
which the recipe's \texttt{/3/tutorial/} filter kept 17.

Where the time went depended on which resource bound the run. USCIS spent
212.7 of its 296 seconds in the streaming stage, close to the 200-second
floor that 20 pages at one request per 10 seconds imposes, and a further
54.2 seconds in discovery fetching its sitemap index and five sub-sitemaps
at the same spacing: the run was bound by the site's declared crawl delay,
as it should be. The Python tutorial, with no crawl delay and nine times
as many chunks as pages, was bound by the LLM instead: 142.5 seconds
streaming and 83.8 seconds in the judge. Ready.gov fell between the two.

The quality stage rejected 20 of 831 samples (2.4\%): 15 for referring to
``the source'' or ``the document'' rather than the subject, 3 for a judge
score below the threshold, and 2 as exact duplicates. None was rejected as
ungrounded and none as a near-duplicate paraphrase. Every export's
training, validation and test files were page-disjoint, as the group-aware
split guarantees, but at this scale the test split of each site held only
one or two pages.

\subsection{Defects the benchmark exposed}
\label{sec:eval-defects}
The unit tests for each component passed before the benchmark was run.
The end-to-end runs nonetheless exposed six defects, all fixed in version
2.4.0, each with a regression test that fails on the previous code:

\begin{enumerate}[leftmargin=1.5em]
  \item \textbf{\texttt{Crawl-delay} was not applied while scraping.}
    DataForge logged ``limiting to 0.100 req/s'' for USCIS, then fetched
    all 20 pages in about four seconds. The flag recording that a domain's
    delay had been applied was process-wide, so the delay reached only the
    first rate limiter created (discovery's); the scraping stage built its
    own limiter and ran at the default rate.
  \item \textbf{The rate limiter ran at about twice its configured rate.}
    The token bucket recorded the time before sleeping, so the next call
    counted the sleep as refill and went through free; it also allowed an
    initial burst of twice the per-second rate.
  \item \textbf{Stage boundaries and jitter shortened individual gaps.}
    With the first two fixed, the median gap to USCIS was 9.99~seconds,
    but the first scrape request followed the last sitemap request by
    3.95~seconds (each stage still had its own limiter), and random jitter
    counted as refill let single gaps fall to 8.6~seconds. After the fix
    (one limiter per run; jitter can only lengthen a gap), the 26 requests
    to USCIS had a minimum gap of 9.73~seconds and a median of
    10.16~seconds, measured between response log timestamps, so variation
    in response latency accounts for the residual shortfall.
  \item \textbf{A shallow sitemap suppressed crawling.} Given the seed
    \texttt{/3/tutorial/}, discovery found the site's sitemap, returned its
    eight unrelated version roots, and dropped the seed; the site could
    not be crawled at all. Discovery now falls back to a crawl from the
    seed when the sitemap lists nothing under the seed's path.
  \item \textbf{Relative links resolved against the site root.} From
    \texttt{/3/tutorial/}, the link \texttt{../genindex.html} was resolved
    to \texttt{/genindex.html} rather than \texttt{/3/genindex.html}, so
    the crawler left the section it was asked to cover.
  \item \textbf{Valid sites failed TLS verification.} Four of the Kenyan
    sites could not be fetched with Python's default certificate bundle: three send an incomplete
    certificate chain and one chains to a root absent from the
    \texttt{certifi} bundle, although browsers accept all four. DataForge
    now verifies against the operating system's trust store; verification
    is never disabled, and self-signed and expired certificates are still
    rejected. (These sites were then excluded on terms, not on access.)
\end{enumerate}

The first three matter most, because Section~2.2 describes
\texttt{Crawl-delay} enforcement as a design property, and it was true of
the component in isolation but not of the assembled pipeline until
version~2.4.0. Politeness guarantees have to be measured end to end, with
timing, against a site that declares them.

\subsection{What these results do not show}
\label{sec:eval-limits}
\begin{itemize}[leftmargin=1.5em]
  \item \textbf{The approval rate is not evidence of quality.} The same
    model generated and judged every sample, and LLM judges are known to
    favour output resembling their own \cite{zheng2023judging}. An
    approval rate of 97.6\% says that the judge rarely objected, not that
    the samples are correct. The repository includes a blind audit
    (\texttt{evals/pipeline/audit.py}) that draws equal numbers of approved
    and rejected samples per site for a person to label without seeing the
    judge's verdict, from which judge precision and rejection precision
    can be computed. Those labels had not been collected at the time of
    writing; we will report them in a later version rather than infer
    quality from the judge alone.
  \item \textbf{The scale is small.} Twenty URLs per site is enough to
    show where time and cost go and to expose the defects above, but not
    to give confidence intervals, and a test split of one or two pages is
    too small to evaluate a model on.
  \item \textbf{URL selection was not random.} \texttt{max\_urls} keeps
    the first N URLs in sitemap order.
  \item \textbf{One model, one run per site, one platform.} Results
    depend on \texttt{gpt-4o-mini}; operating-system trust stores differ,
    and on Linux the store does not fetch missing intermediate
    certificates, so the TLS fix may not help there.
  \item \textbf{No downstream evaluation.} Whether a model fine-tuned on
    this output answers held-out questions better is a separate
    experiment (\texttt{evals/harness/}), not reported here.
\end{itemize}

\section{Agent integration via the Model Context Protocol}
\label{sec:mcp}

The recipe format and distinct exit codes of Section~2 already make
DataForge scriptable. A growing share of that scripting, however, is done
not by a human writing a shell script but by an LLM agent running inside a
coding assistant or chat client. DataForge supports two integration paths
for such agents, at different costs.

\subsection{Shell use with a bundled guide}
The lowest-cost path needs no protocol at all: any agent that can run
shell commands can drive the CLI, provided it knows which commands are
safe to call without a terminal. \texttt{dataforge agent-guide} prints a
guide shipped inside the installed package, so it is available after
\texttt{pip install} without the source repository. It separates the
non-interactive commands (\texttt{run}, \texttt{run --dry-run},
\texttt{init-recipe}, \texttt{explore}, and the \texttt{--json} forms of
\texttt{sessions}, \texttt{stats} and \texttt{view}) from those that open
interactive menus and would hang a non-TTY shell, and states the
configuration, workflow, exit-code and responsible-use rules an agent must
follow. A test asserts that every command the guide names exists, so the
guide cannot drift silently from the CLI.

\subsection{A local MCP server}
The Model Context Protocol (MCP) is an open standard, introduced by
Anthropic in November 2024, for connecting AI applications to external
tools and data sources \cite{anthropic2024mcp}. An MCP server advertises
\emph{tools} with typed input schemas that a client lets its model call.
Hosting is not required: under the protocol's \texttt{stdio} transport,
the client launches the server as a subprocess and exchanges messages over
its standard streams \cite{mcpspec2026}. \texttt{dataforge mcp} is such a
server; in Claude Code, for example, it is registered with
\texttt{claude mcp add dataforge -{}- dataforge mcp}
\cite{claudecodemcp}. The server exposes seven tools
(Table~\ref{tab:mcp-tools}).

\begin{table}[h]
  \centering
  \small
  \caption{Tools exposed by \texttt{dataforge mcp}.}
  \label{tab:mcp-tools}
  \begin{tabular}{@{}l l p{7.6cm}@{}}
    \toprule
    Tool & Effect & Purpose \\
    \midrule
    \texttt{explore\_site}    & Network read   & List a site's sitemap URLs before a recipe is written \\
    \texttt{validate\_recipe} & Read-only      & Dry-run a recipe and return the execution plan \\
    \texttt{start\_run}       & Crawls, spends & Launch \texttt{dataforge run} in the background; return a run ID \\
    \texttt{run\_status}      & Read-only      & Running or finished, exit code and its meaning, session ID, log tail \\
    \texttt{list\_sessions}   & Read-only      & Sessions in the project database \\
    \texttt{session\_stats}   & Read-only      & The session statistics of Section~2.6, as JSON \\
    \texttt{view\_samples}    & Read-only      & Sample records from any pipeline stage \\
    \bottomrule
  \end{tabular}
\end{table}

Three design decisions follow from the pipeline's characteristics rather
than from MCP itself. First, a full run takes minutes to hours --- longer
than a client should hold a single tool call open --- so a run is split
into \texttt{start\_run}, which returns immediately, and
\texttt{run\_status}, which the agent polls. Second, the read-only tools
invoke the CLI's own \texttt{--json} output rather than reimplementing
its queries, so the CLI remains the single source of truth and the two
interfaces cannot disagree. Third, and most important, the safety
constraints that matter most are \emph{enforced by the server} rather
than requested of the model. \texttt{start\_run} refuses a recipe that
sets no spending cap (\texttt{generation.max\_cost\_usd} or
\texttt{generation.max\_llm\_calls}, Section~\ref{sec:budget}) unless the
caller explicitly opts in, and it always refuses a recipe that disables
\texttt{robots.txt} compliance, because an agent has no way to verify the
written permission that setting presupposes; such a run must be started
by the user from their own terminal.

This division is deliberate. The specification requires clients to treat
a server's tool annotations (such as a ``read-only'' hint) as untrusted,
and recommends that a human remain in the loop with the ability to deny
any tool invocation \cite{mcpspec2026}. Those are client-side safeguards.
A rule written only into a model's instructions can be ignored; a rule
enforced by the server cannot. The guide of the previous subsection is
served both as the server's instructions and as a readable resource, so
both integration paths carry the same responsible-use rules
(Section~\ref{sec:responsible}).

\subsection{Potential benefits}
We have not measured the effect of agent access on dataset quality or on
operator effort; the following are expected benefits, not results.
\begin{itemize}[leftmargin=1.5em]
  \item \textbf{Typed interaction instead of text parsing.} Tools accept
    schema-validated arguments and return JSON, so an agent does not have
    to infer structure from terminal output formatted for humans.
  \item \textbf{A lower barrier for domain experts.} An operator can
    describe the target domain in natural language and have the agent
    explore the sitemap, draft a scoped recipe, dry-run it, and report
    the resulting statistics, while the recipe itself remains the
    reviewable, version-controlled artifact of Section~2.
  \item \textbf{Clients without a shell.} Desktop and chat clients that
    cannot run arbitrary commands can still drive the pipeline.
  \item \textbf{Reproducible agent configuration.} Claude Code's
    project-scoped MCP configuration is a file (\texttt{.mcp.json})
    committed to the repository \cite{claudecodemcp}, so the tools an
    agent could use can be versioned alongside the recipes it ran.
  \item \textbf{Composition with other servers.} One agent session can
    combine DataForge with other MCP servers, for example to inspect a
    published dataset or to file an issue about a failed run.
\end{itemize}

The current server keeps its registry of background runs in memory: a
restarted server can still read a run's log but no longer knows its exit
code. It also offers no tool to pause or cancel a run and sends no
progress notifications, so the agent must poll. None of these is
fundamental, but each is a gap in the present implementation.

\section{Related work}

\textbf{Web-scale collection.} Large text corpora for language model
training are commonly built from Common Crawl snapshots with extensive
filtering and deduplication pipelines, as in
CCNet~\cite{wenzek2020ccnet} and the C4 corpus~\cite{raffel2020exploring}.
These systems optimize for scale and general-domain coverage; DataForge
instead targets small, precisely-scoped, single- or few-domain crawls
driven by an explicit recipe, where politeness (rate limits,
\texttt{robots.txt}) and per-run reproducibility matter more than
throughput.

\textbf{Synthetic instruction data.} Self-Instruct~\cite{wang2023selfinstruct}
and its descendants (Alpaca~\cite{taori2023alpaca},
WizardLM~\cite{xu2023wizardlm}) generate instruction-tuning data by
prompting a strong model, typically from a small seed set or from raw
model generations rather than a grounding document. DataForge's
generation stage is closer to retrieval/document-grounded synthetic QA
generation, where each sample is anchored to a specific source chunk ---
which is also what makes the source-reference quality filter and the
group-aware split meaningful, since grounding is both the point of the
dataset and the source of its leakage risk.

\textbf{LLM-as-judge evaluation.} Using a capable LLM to score generated
text has become common for both dataset filtering and model evaluation
\cite{zheng2023judging}. DataForge's quality judge applies this
pattern narrowly --- scoring groundedness of a sample against its own
source chunk, a task closer to a natural-language-inference check than
open-ended quality judgment --- specifically to bound the judge's known
weakness on more subjective, open-ended criteria.

\textbf{Near-duplicate leakage.} Deduplication of near-identical text
across pretraining corpora is known to affect both training efficiency
and evaluation validity \cite{lee2022deduplicating}. DataForge's
group-aware split addresses the same underlying phenomenon at the
dataset-construction stage rather than the corpus-deduplication stage:
instead of removing near-duplicates, it ensures that when they do exist
(as an artifact of $n>1$ samples per chunk), they cannot span the
train/test boundary.

\section{Retrieval augmentation as a complement, not a substitute}
\label{sec:rag}

Fine-tuning on a grounded synthetic dataset and retrieval-augmented
generation (RAG) are often discussed as competing answers to the same
problem --- getting a model to produce accurate, domain-specific output
--- but they address different failure modes and compose rather than
substitute for one another \cite{gao2023ragsurvey}. RAG
\cite{lewis2020rag} retrieves relevant passages at inference time and
conditions generation on them, which bounds hallucination by construction
whenever retrieval succeeds: the model is answering from text placed in
front of it, not from parametric memory alone. What RAG does not change
is the model's underlying style, domain vocabulary, task format, or
ability to follow a narrow instruction schema — that remains a function
of what the model was trained on, which is exactly what a fine-tuning
dataset like the one DataForge produces is for. Conversely, fine-tuning
alone does not eliminate hallucination: a model fine-tuned on grounded
Q\&A pairs learns the \emph{style} of grounded answers, including,
plausibly, the confident tone of a correct one, without any inference-time
guarantee that a novel question is actually answerable from what the
model retained in its weights. Survey evidence on hallucination in
natural language generation \cite{ji2023hallucination} is consistent with
both halves of this: hallucination is not fully solved by better training
data alone, nor by retrieval alone, since a retriever can itself return
an irrelevant or stale passage that a confidently fine-tuned model then
answers from anyway.

A practical implication for a tool like DataForge is that the same
processed corpus it already produces as a byproduct of dataset
construction --- cleaned, boilerplate-stripped, token-bounded chunks with
persisted source lineage (\texttt{page\_id}, \texttt{chunk\_id},
\texttt{source\_url}; Section~2.6) --- is also exactly the unit of text a
RAG retriever indexes. Nothing in the current pipeline exposes this
chunk store as a retrieval index; doing so (e.g., an export target that
writes chunks with embeddings to a vector index rather than, or in
addition to, writing synthetic Q\&A pairs to a fine-tuning-format file)
would let the same crawl produce both a fine-tuning dataset \emph{and}
the retrieval corpus to pair it with, rather than requiring a second,
separate collection effort to stand up RAG over the same domain. We
present this as a natural extension (Section~\ref{sec:limitations})
rather than a shipped capability.

\section{Responsible use and misuse considerations}
\label{sec:responsible}
\label{sec:misuse}

DataForge automates two activities that are each individually subject to
misuse when unattended and applied at scale: web collection, and
synthetic content generation from that collected content. We describe
what the system does to reduce specific, concrete risks, and what remains
the operator's responsibility, since neither an unattended crawler nor an
unattended text generator is safe by default.

\textbf{Collection abuse.} An unattended, resumable, streaming crawler is
easier to run at a scale and cadence that a manual scraping script is
not, which raises the ceiling on how much load an operator can put on a
target site, intentionally or not. DataForge honors \texttt{robots.txt}
by default, including a declared \texttt{Crawl-delay}
(Section~2.3), and exposes \texttt{source.ignore\_robots} as an explicit,
named opt-out rather than a default — a design choice intended to make
disabling politeness a deliberate, auditable act in a committed recipe
file rather than an accidental one. This is a technical control, not a
legal or ethical one: \texttt{robots.txt} compliance does not by itself
establish that scraping a given site is permitted under its Terms of
Service, and DataForge does not attempt to interpret ToS text. The
benchmark made this concrete (Section~\ref{sec:eval-access}): one site
with no \texttt{robots.txt} prohibits data extraction in its terms,
another permits the page in \texttt{robots.txt} but blocks automated
clients, and a third invites all crawlers while reserving all rights. The
benchmark therefore gates each site on a recorded human terms review, and
the agent integration of Section~\ref{sec:mcp} carries the same rule. Rate
limiting defaults to a conservative 2 requests/second/domain; nothing
prevents an operator from configuring a much higher rate against a site
with no declared \texttt{Crawl-delay}; this remains an explicit
operator decision, surfaced in the recipe, rather than a hidden default.

\textbf{Content laundering and provenance loss.} Because generation
produces new text (a question and answer) rather than verbatim excerpts,
the resulting dataset is one step removed from its source in a way that
can obscure where the underlying facts came from, particularly once
exported without its \texttt{source\_url} lineage (e.g., via the
Unsloth-format export path noted in Section~2.6). This is a specific
instance of the broader dataset-provenance problem documented across the
training-data ecosystem, where licensing and attribution information is
frequently lost or misrepresented as data moves between hosting sites and
downstream uses \cite{longpre2023provenance}. Preserving
\texttt{source\_url} end to end, and treating an export format that drops
it as a known gap rather than a cosmetic one, is the concrete mitigation
available today; the ethical and legal obligation to have had permission
to collect and redistribute the underlying content in the first place is
not something the tool can verify on an operator's behalf.

\textbf{Synthetic-data quality claims.} A dataset produced by this
pipeline is only as reliable as its generation goal and quality
thresholds make it, and the tool's own \texttt{quality}-stage terminology
(``approved,'' ``passed'') can read as a stronger correctness guarantee
than a threshold-based LLM judge (Section~2.6, and
\cite{zheng2023judging} on the general limits of LLM-as-judge evaluation)
actually provides, particularly near the score threshold. Presenting an
approved export as an unqualified ground-truth dataset --- rather than as
one scored by a specific, imperfect, and disclosed method --- risks the
same kind of unearned authority that unvetted large-scale text corpora
have been criticized for propagating downstream
\cite{bender2021parrots}. We recommend recipes record and publish their
\texttt{quality.threshold}, \texttt{quality.min\_judge\_score}, and judge
model alongside any released dataset, which the export lineage already
makes possible but does not currently automate.

\textbf{Sensitive and personal content.} DataForge collects whatever a
target site publishes; it has no content classifier that detects or
withholds personally identifiable information, protected health
information, or other sensitive material encountered incidentally during
a crawl. Scoping a recipe's \texttt{source.include}/\texttt{exclude}
patterns away from pages known or likely to contain such content, and
using a local model (Ollama) rather than a cloud provider when source
material carries data-residency constraints, are operator-side
mitigations documented in the project's ethics guidance; they are not
enforced by the pipeline itself.

None of these considerations are unique to DataForge --- they apply to
any tool that automates web collection and LLM-based content generation
at scale --- but an unattended, cron-friendly design (Section~1) is
specifically what makes it practical to run such a pipeline
unsupervised, which raises rather than lowers the importance of stating
these risks explicitly alongside the system's capabilities.

\section{Limitations and future work}
\label{sec:limitations}

The LLM judge is a single model call scoring groundedness on a 1--5
scale; it does not currently ensemble multiple judges or models, which
would likely improve reliability near the threshold at proportionally
higher cost. Its agreement with human judgement is not yet measured: the
blind audit of Section~\ref{sec:eval-limits} is built but its labels are
pending, and in the benchmark the same model generated and judged every
sample. The BFS fallback crawler and SPA (JavaScript-rendered page)
detection are heuristic rather than guaranteed to find all in-scope
content on sites without a sitemap. Group-aware splitting addresses
leakage from the generation stage's own repeated sampling per chunk, but
does not detect near-duplicate content that exists \emph{across} distinct
source pages (e.g., a page and its syndicated copy on another domain);
a content-hash or embedding-similarity based deduplication pass across
groups, rather than only within a session's chunks, is a natural
extension.

Two gaps identified during a post-hoc code review, rather than designed
in from the start, remain open at the time of writing: the Unsloth-format
export's loss of source lineage (Section~2.6), for which the fix is to
carry lineage as a parallel, format-appropriate side channel (e.g. an
adjacent \texttt{dataset\_unsloth.meta.jsonl} keyed by row index) rather
than any change to splitting or generation logic; and the row-order
artifact within a split's output file (Section~\ref{sec:row-shuffle}),
for which the fix is an additional, independent shuffle step rather than
a change to group assignment itself. Both are recorded here as a
reminder that lineage and ordering guarantees need to be verified
per output format and per pipeline stage, not assumed to hold uniformly
once established for one.

The most significant near-term extension, and the one we consider the
project's actual next step rather than a background item, is support for
\textbf{non-text source material}. The pipeline today assumes its input
is HTML prose reachable by an HTTP client: PDFs and images are filtered
out of crawl candidates entirely (Section~2.2) rather than processed,
tabular and structured data embedded in a page is flattened or lost by
the Markdown-conversion step (Section~2.3), and there is no handling for
audio or video sources at all. Domains where the authoritative content is
a table (regulatory filings, statistical releases), a scanned or
image-based PDF (many government and legal documents), or a recorded
briefing are currently out of scope, not because the downstream
pipeline (chunking, generation, quality, leak-aware export) is
text-format-specific in principle, but because Discovery and Collection
only know how to find and fetch HTML today. Extending Collection with a
PDF/table extraction path (text and layout-aware, so a table survives as
structured data rather than becoming word salad) and an
optional OCR or vision-model pass for image-based content, feeding the
same downstream Processing/Generation/Quality/Export stages described in
Section~2, is the direction we intend to pursue next.

\section{Conclusion}

DataForge treats the operational realities of scraping real websites and
the statistical realities of splitting a synthetically-generated dataset
as core design constraints rather than details to fix after the fact.
The result is a pipeline whose default behavior --- honoring server
signals during collection, checkpointing every unit of work, overlapping
collection with generation, bounding unattended LLM spend explicitly,
catching near-duplicate paraphrases rather than only exact matches, and
refusing to let a random split leak near-duplicate content across train
and test --- is also the behavior that produces a defensible dataset and
a trustworthy evaluation number from it. The benchmark of
Section~\ref{sec:eval} adds a caution: three of the six defects it found
concerned politeness, a property that held for the rate limiter in
isolation and failed in the assembled pipeline until it was measured end
to end, against a site that declares a crawl delay. Fine-tuning on that dataset is,
as Section~\ref{sec:rag} discusses, one half of reducing hallucination
rather than the whole of it, and the same corpus the pipeline builds
along the way is positioned to serve the other half (retrieval) as
naturally as it serves fine-tuning. None of this removes the operator's
responsibility for what is collected and published in the first place
(Section~\ref{sec:misuse}); it narrows the space of failures that would
otherwise have to be caught by hand.

\section*{Acknowledgments}
Parts of the DataForge code and of the text of this report were drafted
with the assistance of an AI coding assistant (Claude, by Anthropic), used
under the author's direction. The author reviewed and edited all of it and
takes full responsibility for the content. The figures in
Section~\ref{sec:eval} are taken from the output of the scripts in
\texttt{evals/pipeline/}, which compute them from the recorded benchmark
data, and can be regenerated from it.

\begin{thebibliography}{99}

\bibitem{wei2022finetuned}
J.~Wei, M.~Bosma, V.~Y.~Zhao, et al.
\newblock Finetuned Language Models Are Zero-Shot Learners.
\newblock \emph{International Conference on Learning Representations (ICLR)}, 2022.

\bibitem{ouyang2022training}
L.~Ouyang, J.~Wu, X.~Jiang, et al.
\newblock Training language models to follow instructions with human feedback.
\newblock \emph{Advances in Neural Information Processing Systems (NeurIPS)}, 2022.

\bibitem{wang2023selfinstruct}
Y.~Wang, Y.~Kordi, S.~Mishra, et al.
\newblock Self-Instruct: Aligning Language Models with Self-Generated Instructions.
\newblock \emph{Proceedings of ACL}, 2023.

\bibitem{taori2023alpaca}
R.~Taori, I.~Gulrajani, T.~Zhang, et al.
\newblock Stanford Alpaca: An Instruction-following LLaMA Model.
\newblock \emph{GitHub repository}, 2023.
\newblock \url{https://github.com/tatsu-lab/stanford_alpaca}

\bibitem{xu2023wizardlm}
C.~Xu, Q.~Sun, K.~Zheng, et al.
\newblock WizardLM: Empowering Large Language Models to Follow Complex Instructions.
\newblock \emph{arXiv preprint arXiv:2304.12244}, 2023.

\bibitem{zheng2023judging}
L.~Zheng, W.-L.~Chiang, Y.~Sheng, et al.
\newblock Judging LLM-as-a-Judge with MT-Bench and Chatbot Arena.
\newblock \emph{Advances in Neural Information Processing Systems (NeurIPS)}, 2023.

\bibitem{kaufman2012leakage}
S.~Kaufman, S.~Rosset, C.~Perlich, and O.~Stitelman.
\newblock Leakage in Data Mining: Formulation, Detection, and Avoidance.
\newblock \emph{ACM Transactions on Knowledge Discovery from Data}, 6(4), 2012.

\bibitem{lee2022deduplicating}
K.~Lee, D.~Ippolito, A.~Nystrom, et al.
\newblock Deduplicating Training Data Makes Language Models Better.
\newblock \emph{Proceedings of ACL}, 2022.

\bibitem{pedregosa2011scikit}
F.~Pedregosa, G.~Varoquaux, A.~Gramfort, et al.
\newblock Scikit-learn: Machine Learning in Python.
\newblock \emph{Journal of Machine Learning Research}, 12:2825--2830, 2011.

\bibitem{wenzek2020ccnet}
G.~Wenzek, M.-A.~Lachaux, A.~Conneau, et al.
\newblock CCNet: Extracting High Quality Monolingual Datasets from Web Crawl Data.
\newblock \emph{Proceedings of LREC}, 2020.

\bibitem{raffel2020exploring}
C.~Raffel, N.~Shazeer, A.~Roberts, et al.
\newblock Exploring the Limits of Transfer Learning with a Unified Text-to-Text Transformer.
\newblock \emph{Journal of Machine Learning Research}, 21(140):1--67, 2020.

\bibitem{lewis2020rag}
P.~Lewis, E.~Perez, A.~Piktus, et al.
\newblock Retrieval-Augmented Generation for Knowledge-Intensive NLP Tasks.
\newblock \emph{Advances in Neural Information Processing Systems (NeurIPS)}, 2020.
\newblock arXiv:2005.11401.

\bibitem{gao2023ragsurvey}
Y.~Gao, Y.~Xiong, X.~Gao, et al.
\newblock Retrieval-Augmented Generation for Large Language Models: A Survey.
\newblock \emph{arXiv preprint arXiv:2312.10997}, 2023.

\bibitem{ji2023hallucination}
Z.~Ji, N.~Lee, R.~Frieske, et al.
\newblock Survey of Hallucination in Natural Language Generation.
\newblock \emph{ACM Computing Surveys}, 55(12), Article 248, 2023.

\bibitem{bender2021parrots}
E.~M.~Bender, T.~Gebru, A.~McMillan-Major, and S.~Shmitchell.
\newblock On the Dangers of Stochastic Parrots: Can Language Models Be Too Big?
\newblock \emph{Proceedings of the 2021 ACM Conference on Fairness, Accountability, and Transparency (FAccT)}, pages 610--623, 2021.

\bibitem{longpre2023provenance}
S.~Longpre, R.~Mahari, A.~Chen, et al.
\newblock The Data Provenance Initiative: A Large Scale Audit of Dataset Licensing \& Attribution in AI.
\newblock \emph{arXiv preprint arXiv:2310.16787}, 2023.

\bibitem{anthropic2024mcp}
Anthropic.
\newblock Introducing the Model Context Protocol.
\newblock \url{https://www.anthropic.com/news/model-context-protocol}, November 25, 2024.

\bibitem{mcpspec2026}
Model Context Protocol.
\newblock Specification, revision 2026-07-28 (Transports; Server Features: Tools).
\newblock \url{https://modelcontextprotocol.io/specification/2026-07-28}, accessed 2026-09-22.

\bibitem{claudecodemcp}
Anthropic.
\newblock Connect Claude Code to tools via MCP.
\newblock \url{https://code.claude.com/docs/en/mcp}, accessed 2026-09-22.

\bibitem{usc17sec105}
United States Code, Title 17, Section 105.
\newblock Subject matter of copyright: United States Government works.

\bibitem{kenyacopyright2001}
Republic of Kenya.
\newblock Copyright Act, No.~12 of 2001 (Cap.~130), Revised Edition 2022.
\newblock \url{https://copyright.go.ke/sites/default/files/downloads/CopyrightAct12of2001.pdf}, accessed 2026-09-22.

\end{thebibliography}

\end{document}
